\section{Human Evaluation Study}
\label{sec:human_study}

We conduct a human evaluation study to validate LLAMIA's commentary generation against automatic metrics. Unlike subjective preference studies, we design task-based evaluations that measure whether commentary enables accurate game understanding.


\subsection{Participants and Protocol}
\label{sec:study_participants}


We recruit 12 participants from two institutions with \url{chess.com} affiliated programs. All participants meet the following criteria: (1) peak online rating of 1600+ on chess.com or lichess.org, (2) over 20,000 online games played, and (3) prior annotation experience on GameKnot.com's game analysis platform. Participants were presented with a study protocol and signed IRB approval from the participating institution, which detailed the data collection procedures and usage terms.


\subsection{Commentary Evaluation Tasks}
\label{sec:commentary_tasks}


We compare commentary from three sources: Human expert commentary (from the Agadmator dataset), LLAMIA-13B, and GPT-4o with Lc0 tool access. For each of 50 game segments (15--25 moves), participants complete three tasks designed to measure functional understanding rather than subjective preference.

\paragraph{Task 1: Position Evaluation from Commentary (Accuracy).}
We test whether commentary accurately conveys positional assessment. At 5 randomly selected positions within each game segment, participants read only the commentary (no board visible) and predict the engine evaluation on a 7-point scale: decisive White advantage ($> +3$), clear White advantage ($+1.5$ to $+3$), slight White advantage ($+0.5$ to $+1.5$), equal ($-0.5$ to $+0.5$), slight Black advantage, clear Black advantage, decisive Black advantage. We measure RMSE between predicted and actual centipawn evaluations (scaled to the 7-point range).

\paragraph{Task 2: Move Prediction from Commentary (Insight).}
We test whether commentary explains strategic rationale well enough to guide move selection. At 5 positions per segment, participants read the commentary describing the current position and are presented with three legal move options (the played move plus two engine-suggested alternatives). Participants predict which move was played based on the commentary's description of plans and threats. We report prediction accuracy.

\paragraph{Task 3: Quality (Engagement).}
Participants rate each commentary segment on a 5-point Likert scale: ``This commentary would help a 1400-rated player understand the strategic themes of this game.'' This captures perceived pedagogical value.


\subsection{Results}
\label{sec:study_results}


\begin{table}[h]
\centering
\begin{adjustbox}{max width=0.98\textwidth}
\begin{tabular}{l|cc|c}
\toprule[1.2pt]
\textbf{Commentary Source} & \textbf{Eval Prediction} & \textbf{Move Prediction} & \textbf{Educational Value} \\
 & (RMSE $\downarrow$) & (Accuracy $\uparrow$) & (Likert 1-5 $\uparrow$) \\
\midrule
Human (Agadmator) & \valbest{0.89} & \valbest{71.2\%} & \valbest{4.31} \\
LLAMIA-13B & \valgood{1.12} & \valgood{64.8\%} & \valgood{3.89} \\
GPT-4o + Lc0 & 1.47 & 58.3\% & 3.42 \\
\midrule
Random Baseline & 2.31 & 33.3\% & -- \\
\bottomrule[1.2pt]
\end{tabular}
\end{adjustbox}
\caption{\textbf{Commentary Evaluation Results.} LLAMIA-13B achieves lower evaluation prediction error (1.12 vs 1.47 RMSE) and higher move prediction accuracy (64.8\% vs 58.3\%) than GPT-4o + Lc0, approaching human expert performance (0.89 RMSE, 71.2\% accuracy). Educational value ratings follow the same pattern (3.89 vs 3.42 Likert). Results are averaged over 12 participants and 50 game segments (250 evaluation points per commentary source). Inter-rater reliability: Krippendorff's $\alpha = 0.74$ for evaluation prediction, $\alpha = 0.81$ for move prediction.}
\label{tab:commentary_study}
\end{table}

\paragraph{Analysis.}
LLAMIA's advantage over GPT-4o + Lc0 is most pronounced on evaluation prediction (24\% lower RMSE), indicating that latent-grounded commentary better conveys positional assessment than tool-calling approaches. The gap narrows on move prediction (6.5\% higher accuracy), suggesting both systems adequately describe immediate tactical considerations. Human commentary remains superior, particularly on evaluation prediction (21\% lower RMSE than LLAMIA), reflecting experts' ability to synthesize complex positional factors into coherent narratives.

\paragraph{Qualitative Patterns.}
Participants noted that LLAMIA commentary more frequently mentions specific positional features (``weak d5 square'', ``backward e-pawn'') compared to GPT-4o + Lc0, which tends toward generic assessments (``White has a slight advantage''). However, human commentary excels at narrative continuity, connecting individual moves to longer-term strategic arcs.